\section{The Hidden Tradition}
\label{sec:HiddenTradition}

\begin{quotex}
For most men life is made up of two days:

\begin{enumerate}
\item 
In the first they believe everything. 
\item
And in the second, nothing. 
\end{enumerate}

For a few others, life also has two days, but what distinguishes them from ordinary men is that

\begin{enumerate}
\item
In the first day they believe only in illusions, and these are nothing; 
\item
While in the second day, they believe in everything, for they believe in truth, which is all. 
\end{enumerate}
\flright{\textsc{Louis-Claude de Saint-Martin}}

\end{quotex}
Since \textbf{Rene Guenon} was an initiate into a \textbf{Martinism} order, it may be of interest to explore their teachings. Martinism derives its name from the two ``Martins": \textbf{Martinez Pasquales} and \textbf{Louis-Claude de Saint-Martin}. The former mysteriously appeared with an order based on theurgy or ceremonial magic. The latter joined that order, but eventually abandoned it when he felt that he moved beyond ``operations". There was a loose group that gathered around him, the S.I. for the Order of Superieurs Inconnus (Unknown Superiors) and initiation was decidedly less formal.

Now Martinism claims its origin from God which was communicated from Adam to Noah to Melchisedek and then Abraham, Moses, Solomon, and Zerubbabel. In other words, there has always been this hidden tradition. \textbf{Valentin Tomberg} called it the Church of John. In a sense Tomberg completes this teaching, since he claims there is no longer a need for an ``order". Furthermore, the uneasy, and often hostile, relationship between the Johannine and Petrine churches is to be rectified.

\paragraph{Martinist Doctrine}
The doctrine hinges on a particular view of the fall of man; these are the fundamental points:

\begin{enumerate}
\item Original Man (or Primitive Adam or Archetypal Man) was emanated by God and dwelt on a high plane where he enjoyed a privileged position. 
\item Through the misuse of his free will, Primitive Adam sinned and as a result fell. 
\item As a result of the Fall, Man, originally a unified being, became shattered into numerous individualities who are now the men and women of the material world. 
\item Individual man's task is to reintegrate himself back to the Archetype. All men and women must reintegrate with the Archetype so that Primitive Adam can be reconstituted and unity be again achieved. 
\end{enumerate}
The purpose of the order was to provide a way to achieve that task of reintegration.

\paragraph{Gregory of Nyssa on Adam}
Lest anyone think that is absurd or even heretical, I refer you to \textbf{Gregory of Nyssa}, as explained by \textbf{Vladimir Lossky} in \emph{The Mystical Theology of the Eastern Church}:

\begin{quotex}
The person, however, transcends and is free from his nature, i.e., not determined by it. The individual, on the other hand, which is the material manifestation and individuation, is often confused with the person. Thus, it is the individual which must be transcended. The Person is in a formless state. After the Fall, human nature became divided and broken into many individuals. Man has a double character: as an individual nature, he is a part of a whole, but as a person, he is in no sense a part, he contains all in himself. 

\end{quotex}
So in our language the individual is a lower state of being and the person is a higher state. Human nature is disintegrated into multiplicity; in theological terms, human nature lost the ``Integrity" which Adam possessed. It is man as individual that is transcended, not man as person. In the latter state, he is in the image of God and able to reach the Divine Union.

\paragraph{Saint-Martin's Principles}
Saint-Martin was attracted to the spiritual and esoteric life from an early age and remained a lifelong Catholic. He tells us the principle he discovered that guided his life:

\begin{enumerate}
\item Despite all the confusions of philosophy, he had attained certitude as to God and his own soul. 
\item The seeker for wisdom had need of nothing more. 
\item The foundation of happiness lies in contentment with the truth. 
\item Absorption in material things was incomprehensible for those who knew the treasures of reason and the spirit. 
\item Human science explained matter by matter, but after its putative proofs, there were other demonstrations needed. 
\item The inmost prayer of his soul was for God to abide therein to the exclusion of all else. Thus Divine Union is the true end of man. 
\item We are all widowed and are called to a second marriage. 
\end{enumerate}
Nothing could be clearer. He starts with knowledge of the self and knowledge of God. He ends with the Divine Union, which he relates to a type of marriage. In brief, we can refer to a couple of high points in his doctrine, specifically as they relate to recent discussions.

\paragraph{The Will and Silencing the Mind}
Saint-Martin makes clear that the Divine Union can be accomplished in this life, rather than postmortem. First of all, a man must become aware of himself as a Person. He writes:

\begin{quotex}
When we have once sensed our soul, we are left in no doubt as to its possibility. 

\end{quotex}
The individual whose inner content is filled with opinions, desires, emotions, likes, dislikes, and sensations will experience nothing but doubt. To overcome that, the will must be conformed. He tells us:

\begin{quotex}
Let me affirm that divine union is a work which can be accomplished only by the strong and constant resolution of those who desire it; that there is no other means to this end but the persevering use of a pure will, aided by the works and practice of every virtue, fertilized by prayer, that divine grace may come to help our weakness and lead us to the term of our regeneration. 

\end{quotex}
The Martinist way now makes sense. It is not a question of overcoming human nature, but of transforming the fallen nature into something higher. Without understanding the state of Fallen man, one can only interpret such writings as overcoming human nature itself. That is not so; rather the goal is to transcend the fallen state and rise back to the Primordial State (``reintegration"). Clearly, man in the fallen state cannot accomplish that on his own, but requires an impetus from above him. This is not an indication of passivity, since it is the Will, the very source of activity, that must be employed.

To recap thus far: first, a man becomes aware of the possibilities of his soul, which makes him dissatisfied with ordinary life. Then he purifies the Will, which is no longer interested in ordinary pursuits. However, he still does not know how to proceed without an impetus from above. Unfortunately, the lot of mankind, although steeped in ignorance, is to make stuff up. Saint-Martin writes:

\begin{quotex}
It is a sad vision indeed, that when we start looking at Man, we see him at the same time tormented by the desire to know, failing to find a reason for anything, but daring nevertheless to give a reason to everything. 

\end{quotex}
Hence, man must first admit that he knows nothing, so he can follow his path open to Divine Grace. Otherwise, the knowledge he imagines he possesses will form an insurmountable obstacle to his progress.

\paragraph{The Active Path}
Although Saint-Martin uses the word ``mysticism", he technique is more active. First of all, let's see how he understand mysticism:

\begin{quotex}
The only initiation which I preach and seek with all the ardour of my soul is that by which we may enter into the heart of God and make God's heart enter into us, there to from an indissoluble marriage … 

\end{quotex}
Once again, there is that alchemical imagery with the idea of ``work". The marriage of the soul with God is both a union and a separation. That is the proper meaning of ``non-duality" … not two, but not one either. A marriage is the apt symbol, since it is a union without annihilating either one. For ``God" to annihilate the man is a form of monism, not non-dualism.

The Inner Way of Saint-Martin is individual and particular for each person. Concentration and contemplation are actually active, not passive. Try it, and you will see how difficult it is to maintain the focus of the mind. In the book which we shall mention below, his way is compared to Saint John of the Cross, who says that contemplation is infused passively so that illumination is awaited rather than sought. Saint-Martin sought to establish the correspondence of the soul with the Divine by the active path of WORKS. For the path of the Inner Way, we

\begin{quotex}
seek to explain material things by man, and not man by material things. 

\end{quotex}
That certainly deserves more thought than the few seconds required to read it. Who is prepared to understand the entire course of his life and the manifestation of his inner states?


\hfill

\paragraph{Acknowledgment}
For the source of the doctrines and quotations, I relied on the little book, Five Christian Principals, by ``R.C.", which describes the spiritual sources of the Martinist Doctrine. What follows is my own interpretation of their significance.

\begin{enumerate}
\item Martinez Pasquales. The link to the hidden tradition of all ages. 
\item Louis-Claude de Saint-Martin. The link from the hidden tradition to the Roman church. 
\item Jacob Boehme. The culmination of the current of German mysticism that began with Meister Eckhart. 
\item Thomas a Kempis. Another German, he provides the basis of the path as the imitation of Christ. 
\item Emanuel Swedenborg. He restores the tradition of the symbolic interpretation of Scriptures. 
\end{enumerate}


\flrightit{Posted on 2013-05-20 by Cologero }

\begin{center}* * *\end{center}

\begin{footnotesize}\begin{sffamily}

\texttt{Caleb Cooper on 2013-05-21 at 09:03 said: }

It's worth noting that in 1917 Valentin Tomberg was initiated into Martinism at the age of 17, Providently right before the Communist Revolution convulsed Russia and forced his family to flee for their lives to the West. 
\url{http://en.wikipedia.org/wiki/Valentin\_Tomberg}


\hfill

\texttt{Avery Morrow on 2013-05-21 at 11:06 said: }

Is Swedenborg part of Martinism? He's certainly an integral part of the most ridiculous kinds of New Age spiritualism where all metaphysics goes out the window.


\hfill

\texttt{Jason-Adam on 2013-05-21 at 13:58 said: }

I have issues with this – if Saint-Martin was so Catholic why did he refuse the sacrament of extreme unction at his death ? 

According to Abbe Barruel, Memoirs sur l'histoire du Jacobinisme, Tome 2, Martinism is egalitarian and a predecessor to the French Revolution.

Also, according to Barruel and Leon de Poncins, Martinez was a Jew. Could his ``tradition" be seen as the introduction of false Judaic-derived cabalistic (countertradition) poison into Catholic France ?

Joseph de Maistre, speaking as the character of Le Comte in Dialogue 11 of Les soirees du Saint-Petersbourg refutes martismism (considered as part of the illuminist movement) and defends the Church.


\hfill

\texttt{Cologero on 2013-05-21 at 23:23 said: }

Perhaps, Jason-Adam, Joseph de Maistre is an ambiguous source since he was an initiate into Martinism and knew Saint-Martin personally. Charles M Lombard's intellectual biography of de Maistre, although it border on hagiography, documents the subtle ideas of Martinism in de Maistre's works. It may be a good project someday to verify his references. I would see de Maistre's cautions about ``illuminism" to be like Tomberg's, viz., it needs to be brought back into the fold, so to speak.

The point, however, is that Saint-Martin fulfills pretty well Evola's requirements for a Christian initiation. That will require a decision: to follow or to reject that path. Saint-Martin probably rejected last rites because he felt he achieved the ``true immortality" that Evola described.

As for de Maistre's defence of the Church, who today is willing to defend him? Apart from Cioran and Gornahoor, that is. Stendahl complained that the Jesuits who trained the French priests were imbued with de Maistre, and therefore called him the ``French St Paul". That would be inconceivable in our time.


\hfill

\texttt{Constantine on 2013-05-22 at 09:18 said: }

I wonder would Evola or Guenon approve Mormonism?


\hfill

\texttt{Cologero on 2013-05-22 at 12:30 said: }

Guenon wrote a rather longish essay contra Mormonism, included in ``Miscellanea". What in it warrants approval?


\hfill

\texttt{Jason-Adam on 2013-05-22 at 13:17 said: }

Plinio Correa de Oliveira wrote, the counterrevolution is taking place against the subversive elements within the Church. It is a hard fight we may not win but we fight for truth and honour, not for vain material victories. That is what it means to a knight. ``Saint-Martin fulfills pretty well Evola's requirements fora Christian initiation" – but Evola believed in hierarchy, are you saying that the Abbe Barruel was wrong to classify Saint-Martin as an egalitarian and a predecessor of the French Revolution ?


\hfill

\texttt{Mercurius on 2013-05-22 at 14:19 said: }

It is difficult to understand what Mr. Adams means by ``false Judaic-derived Cabalistic derived counter tradition poison"–does this imply that Cabala is ``counter-traditional" (not just anti-traditional either, but counter-traditional) as an esoteric doctrine, or is it counter-traditional because it is ``Judaic derived"? In either case, Cabala certainly entered Christianity long before Martinez–examples can be found in primitive Christianity, and St. Jerome himself, who gives the Vulgate was considerably influenced by Cabala as seen in his letter to Marcella discussing the Divine Names (which are the names ascribed to the Sephiroth). 

The question of whether Martinism ``is" or ``isn't" egalitarian is a bad question. If something represents an initiatic doctrine (and method), then by its nature, spiritual realization/the spiritual order is its sole concern–if, as Guenon pointed out, an organization having initiatic leanings at some point elevates social and political concerns to a rank of primary importance, it is no longer initiatic de jure, but a ``secret society". So, to know what Martinism is one ought look first to its major ``landmarks"–which can be said to be Martinez ``Treatise on the Reintegration of Beings", the ``degrees" which became grafted onto an Masonic structure (so, the ``form" they took for expression), and then, the writings of St. Martin, either as addendum, continuation, completion, or expression not of theurgical operations, but the ``path of the heart" (it is interesting that a sort of similar thing can be observed in neo-Platonism, with say Proclus' school stressing contemplation, and representatives like Iamblichus arguing for theurgy; or the shift from Vedic ritualism toward Vedanta and the beginnings of yoga–but I digress). At any rate, there isn't anything in those landmarks pertaining to a social/political ``egalitarianism"–consequently, the question doesn't have much meaning when thinking about Martinism as initiatic doctrine and method.

Now, on the other hand, if in time, certain lodges and individuals within Martinism have been, or promoted egalitarianism, then that was/is their proclivity, secondary or tertiary ideas caused to become associated with their initiatic work, basically innovations. Interestingly, in his book ``Beyond Enlightenment", Harvey shows many arguments supporting the thesis that Martinism was anti-semitic, racialist, and proto-fascistic! But again, if some lodges could be found as such–neither egalitarianism nor ``proto-fascisim" constitute the landmarks represented by Martinez ``Treatise" simply because it never was set forth as a social-political doctrine. ``Bringing it back into the fold" as Cologero notes above.

Some of the criticisms above seem to be hinged around an ``us versus them" criteria–or the (Roman) Church versus Martinism–but this too is a non-issue, as such a conflict doesn't exist. Martinism is not a religion, and never had desired to replace the Church–Martinez and St. Martin were Catholics, and did not repudiate exoteric religion; rather, they attempted to ``complete" it with a doctrine/method suitable for the qualified (for the non-qualified, the Church always still remains).


\hfill

\texttt{Jason-Adam on 2013-05-22 at 18:14 said: }

1. The Cabala is not Christian so therefore any system based on it is not a Christian initiation, am I wrong or right ?

2. Many writers have shown the anti-Western nature of Judaism, even Evola, a better writer in this regard is Leon de Poncins. 

3. There are many supposed initiatic organisations that are egalitarian, Franc-Masonry to name one……..one area Evola was superior to Guenon was in his recognising Masonry as a false, evil group………

4. I am relying on the Abbe Barruel, Memoirs sur l'histoire du Jacobinisme, Tome 2, for my information on the egalitarian and revolutionary nature of Martinism. If anyone would care to read my source and offer information that refutes it, I am willing to change my opinion.


\hfill

\texttt{Ash on 2013-05-22 at 19:27 said: }

The writings from the eastern Church are rather apt, as Martin's doctrines appear to correlate very closely with the Orthodox doctrine of theosis. However, as I understand it, the Eastern Church goes beyond restoring man to the primordial state and intends to set man back on the path which they were on in Eden, suggesting that Adam in Eden has himself not yet attained his fullness; perhaps Adam is not yet Adam Kadmon? C. S. Lewis did a marvelous job exploring this idea in ``Perelandra". Perhaps someone more informed on Martinist doctrine could elaborate on how that could affect Martinism as a traditional path.


\hfill

\texttt{Francis Mercuri on 2013-05-22 at 19:43 said: }

Mr. Adam:

1. I suppose this entirely depends on how you define Christianity–my perspective understands it as being a union of the Hebraic and Hellenistic traditions. Yet, even if one were to disregard the Hellenistic elements, being Hebraic (at least in terms of having an OT relationship or foundation), one could certainly see how Hebrew esotericism would, or could be, part of Christian theosophy–I don't think that is stretching the logic too far. On the other hand, I'm not so certain we share the same idea concerning what Cabala happens to be.

Still, looking to our Traditional authors, in Schuon we find him several times insisting that there can never be a complete understanding of scriptures of any type in translation(``Gnosis is not Just Anything", ``Some Difficulties Found in Sacred Scripture")–one must, he argues, always return to the sacred tongue–in terms of the Torah, naturally Hebrew. Then, in the essay ``Keys to the Bible", he goes even further, saying:

``In order to understand the nature of the Bible and its meaning, it is essential to have recourse to the ideas of both symbolism and revelation….for the literal meaning of the Bible….affords sufficient food for piety apart from any question of symbolism, but this nourishment would loose all vitality and all its liberating power without an adequate idea of revelation or of superhuman origin….When approaching scripture, one should always pay the greatest attention to rabbinical and cabalistic commentaries–and in Christianity–to the patristic and mystical commentaries"

2. Evola wrote many things about ``Judaism", often times in apparent contradiction–in Men Among the Ruins, he seems to say that he uses the word ``Jew" and ``Jewish" to denote an economic Judaism (and he draws much from Werner Sombart on all this), or ``spirit"–akin to what he writes about the ``races"–so, it pertains in certain contexts more to a ``type", because in other places, such as ``Revolt", he has positive things to say about Cabala, and quotes from the Zohar, with the particular quote representing he states, the very meaning of the traditional world. It isn't so easy to be cut and dry with such things. 

I tend to agree with the Gornahoor position on the issue, as expressed in the post on ``The Nine Worthies":

``The characters are also archetypal and the symmetry of the scheme reflects the Medieval view of history and destiny of European man. The Old Law prepared the way for the New. The Pagan law created the Pax Romana that allowed the spread of the New Law. And of course, medievals saw themselves as the true heirs of both the Hebrews and the Pagans, and their civilization as the fulfillment of God's divine plan".

http://www.gornahoor.net/?p=331

3. If you considered what I previously wrote, i.e., there is a difference between an Initiatic Order, and a ``secret society", you would see that the comment is null, because if Francmasonerie can be regarded as ``egalitarian", one is speaking not of an Initiatic doctrine and method or a metaphysics, but of a political-social construct, which is secondary. When it comes to Masonry, that is where I think Evola was often at his worst–for he could not see it beyond what had emerged as political Masonry, he could not, as the ``conspiracy theorists" of today continue, disentangle an acquired Jacobinism from an Initiatic process and symbolism in Masonry–something Evola hardly touched upon, but is all through Guenon, and heavily in ``Great Triad", ``Esotericism of Dante", and of course his writings assembled into ``Studies on Masonry and Compagnonnage.

4. One could say similar things about the Daoist ``Triads"–with no, or minimal ``Jewish" influences involved, Triads have historically gotten involved in ``revolutionary" activities–often to the point of eliciting governmental crackdowns; yet, again, one can't simply take the political activities of an isolated Triad as descriptive of Daoism–that is just kind of common sense. We are likely getting far afield though–so, since the essay ``Hidden Tradition" pertains foremost to the doctrines of Martinism, I'd suggest our discussion center more on the Martinist treatises and symbolism–allowing them to do the speaking about themselves.


\hfill

\texttt{Cologero on 2013-05-22 at 22:54 said: }

In the Spiritist Fallacy, Guenon points out that there are many nuggets in Swedenborg if one knows how to find them. However, he is to be interpreted symbolically, not in a new age fashion. Nevertheless, Guenon concludes that what is good in Swedenborg can be found elsewhere in a purer form.


\hfill

\texttt{Mihai on 2013-05-23 at 03:29 said: }

``my perspective understands it as being a union of the Hebraic and Hellenistic traditions. "

To regard Christianity as nothing but this union of these two traditions is to view it the same way as the profane historians do.

Certainly, Christianity assimilated and made much use of Hellenistic thought, but the latter is the support, while the former the essence. There are essential points in which Christian doctrine clearly surpasses or fulfills the thought of a Plato, Aristotel or the neo-platonists. The failure to realize this is what caused the scholastic collapse. The same goes, I believe, more clearly in relation to Judaism.

``When it comes to Masonry, that is where I think Evola was often at his worst"

Evola considered the implications of political and subversive masonry, which, I believe, is all that truly remains of Masonry in the modern world. If there are any traditional organizations left in that field, they are probably so well hidden and reclusive that they are close irrelevant for all purposes regarding initiation and spirituality for those seeking them.

I have always stated that Guenon's defense of Masonry, beyond showing what used to be traditional in the organization, marks the lowest point of his writings. For much of the time while he was in Europe he struggled to make Masonry act as an esoteric organization in order to fulfill what he thought was the entirely ``esoteric-less" Christianity. Not only was such a point fruitless, but it may really fuel some of the subversive forces he so much tried to counteract. 

He had really made better use of his time to find in the Patristic mystical writings the starting point for the reconstruction of a higher perspective for the West, instead of wasting time trying to revive a long dead and buried organization.


\hfill

\texttt{V.O. on 2013-05-23 at 10:46 said: }

Christ annihiliated the boundary between the esoteric and the exoteric. The only ``initiation" now occurs exoterically in the Orthodox Church. Mysticism as the ``estoeric" is practiced by forms of the anchorite, perhaps away from the world, but still within the Orthodox Church. (Although the Orthodox Christian is always away from the world)

The ascetic alone with the wolves in the mountains is exoteric because the Church is universal and the Cross transforms the world. Mysticism as term is a terminological advancement away from the exoteric-esoteric split. The initiate as the preserver of the esoteric and exoteric split is Antichrist, because their absolute gap is not between God and the world, but within the world itself. They are only concerned with the world. The ``secret" organizations are thus always practicing a form of counter-initiation and inevitably degrade into complete Satanism.


\hfill

\texttt{francismercuri on 2013-05-23 at 12:06 said: }

It was never suggested that Christianity was ``just"–or ``merely", anything–as if an external assembly of lifeless fragments from ``cultures". One would assume that as understood though in such a forum as this, hopefully without needing to emphasize the obvious repeatedly. I'll try to be more lucid hereafter.

So called political and ``subversive" Masonry are certainly not (and this can be looked at entirely objectively) ``all that remains". What is there in regular Masonry are all the same landmarks, rites, symbols, and lectures that have been there since the operative went speculative–the very same things that Guenon would have worked with (granted, there are some variations between Grand Lodges in wording, and ``floor work" due to pedestrian and insignificant details)–there is no need for anything to be ``hidden and reclusive". Everything is entirely present if one desires a ``virtual" initiation, and wants to receive such in the context of an ``order". All this back and forth stuff about masonry is silly–if what is contained in Masonry in terms of the rites and symbols is ``political and subversive", then the bulk of ``esotericists" ought throw in the towel, because these are the same symbols and doctrines that are present in other forms of Western esotericism–say Hermeticism/alchemy for example. 

Aside from all that, all this talk of political masonry is in day-to-day reality something of a tremendous exaggeration–if not completely false! ``Political" lodges are likely the exception, if you can even find one–such topics are not even permitted in lodges, they just are not discussed. What lodges do is this: they occupy themselves with performing degree work (initiations)–this takes up 75\% of members time, between practice, memorization, and performance–after which, it is up to the initiate to either penetrate the deeper meanings of the rites–or, do nothing at all–only ``research lodges" with a ``philosophical" bent talk about this stuff. The remaining time, members sometimes spend socially, relaxing, outings, and/or doing some charity work–hospitals, disaster relief–whatever. That is the true life a regular lodge–all these other things people talk about are fantasy.


\hfill

\texttt{Jason-Adam on 2013-05-23 at 21:34 said: }

I fully agree with you !!!!!

Eliphas Levi had said something similar in his History of Magic and so did Fritjof Schuon. 

not for nuttin but the first 2 sacraments are known as the sacraments of INITIATION……


\hfill

\texttt{Jason-Adam on 2013-05-23 at 21:42 said: }

Louis-Claude de Saint-Martin was a supporter of the French Revolution, I can provde sources for this if needed. How can you be a reactionay and support this doctrine then ?

Guenon ABANDONED Martinism when he realised its counterfit nature.

Mr Mercuri, how do you know that stories of political Masonry are exagerrated ? are you a Mason ? if you're not, you can't know for sure because Masons are sworn not to reveal what happens in the lodge to non-Masons. If you are a Mason then you may be trying to cover up your activities…..

I choose to trust the irrefutable testimony of many researchers who know the evils of Satanic Masonry…..

My interest in esotericism these days tends toward the Hieron du Val d'Or…….Cologero have you ever read about that group ? Guenon was associated with them for a time.


\hfill

\texttt{Cologero on 2013-05-24 at 01:09 said: }

Jason-Adam, I've never located any primary literature from the Hieron; Guenon's involvement stems from his association with Regnabit and the correspondence between Guenon and Louis Charbonneau-Lassay. I don't know if that has been published but let me know if you can find it.

I've never found any writings from the founders of the Hieron, so I can't say if the description on Cesnur is accurate. Nevertheless, some of the themes have been take up here. E.g., the common tradition, which has been mentioned several times. Another is the Christianization of Masonry; this, we believe, was ultimately the task of Valentin Tomberg. Apparently, that group was instrumental in promoting the feast day of Christ the King. I will write on that soon provided I can make it of general interest. I think it promotes a better societal vision than what any of the new pseudo-right groups have been able to do. The recent meaningless gesture at Notre Dame has made this suddenly a relevant topic.


\hfill

\texttt{Mihai on 2013-05-24 at 03:35 said: }

Granted, when it comes to certain so-called conspiracy theories (although I really detest the term, since it is used by those who deny any hidden activities that influence historical events) there are many exaggerations.

But there is an old saying, which I strictly keep in mind when it comes to things such as these: where there's smoke, there's fire. It is a fact that Masonry has participated in almost every revolutionary movement of the last few centuries- including the American revolution, the French revolution, the 1848 movements and even the Soviet Revolution. If these days Masonry has ceased to be so political is because there is no further need for it as a platform to spread certain subversive goals, since there are other, more adequate ones.

Anyhow, the fact that it was possible, back then, to use Masonry for subversive purposes, denotes that it had already degenerated beyond any relevance. 

``Everything is entirely present if one desires a ``virtual" initiation, and wants to receive such in the context of an ``order". "

My question is: even if such were true, and there is still some spiritual influence acting in those lodges to pass a valid initiation , I don't see what good it could do if its members do not distinguish themselves, in their mentality, from the least of the profanes and can provide no significant guidance or knowledge to anyone.


\hfill

\texttt{Jason-Adam on 2013-05-24 at 17:37 said: }

I've had the same problem in detecting primary sources as well, Charbonneau-Lassay interests me greatly though, his train of thought today seems to be represented by Jean Borella whom I can recommend strongly. 

As for the monsieur who was a suicide, I believe that is a representation of utter nihilism. For myself as a knight, I do not care what the world around me is like – it is sinful always. I follow God and do not depend on others for my inner serenity. 

I did some more research and Eliphas Levi did admire Saint-Martin and said there was a revolt of the subversive illuminists (Weishaput and co.) against the legitimate authority, Barruel perhaps was too exoterically minded by lumping all the esotericisms into the category of egalitarian and atheist. I dont know enough yet. I still distrust Martinism due to de Maistre and Guenon's rejection of it but I keep my mind open to the future work Gornahoor will undertake on it.


\end{sffamily}\end{footnotesize}
